\documentclass[11pt]{article}
\usepackage[paperwidth=210mm,paperheight=5000mm,margin=20mm]{geometry}
\usepackage{amsmath,amssymb,xcolor,graphicx}
\usepackage[hypcap=false]{caption}
\usepackage[hidelinks,bookmarksnumbered,bookmarksopen]{hyperref}
\hypersetup{pdftitle={Seminar — Garnet solid electrolytes for lithium metal batteries},pdfauthor={Speaker: Dr P. Varga},
  pdfcreator={KherveNote}}
\renewcommand\labelitemii{\textbullet}
\renewcommand\labelitemiii{\textbullet}
\renewcommand\labelitemiv{\textbullet}
\renewcommand\labelenumii{\arabic{enumi}.\arabic{enumii}.}
\renewcommand\labelenumiii{\arabic{enumi}.\arabic{enumii}.\arabic{enumiii}.}
\renewcommand\labelenumiv{\arabic{enumi}.\arabic{enumii}.\arabic{enumiii}.\arabic{enumiv}.}
\setlength{\parindent}{0pt}
\setlength{\parskip}{0.6em}
\definecolor{knoteaccent}{HTML}{1A6DD8}
\definecolor{knotemuted}{HTML}{6B6B6B}
\definecolor{knotekey}{HTML}{D96B00}
\newcommand\knotetime[1]{\leavevmode\llap{\color{knotemuted}\scriptsize #1\hspace{1em}}}
\newcommand\knotekey[2][]{\par\noindent#1\colorbox{knotekey!12}{\parbox{\dimexpr\linewidth-2\fboxsep}{%
  \textbf{\color{knotekey}$\star$ Key point.} #2}}\par}
\newcommand\knotequestion[2][]{\par\noindent#1{\color{knoteaccent}\textbf{?}}~\emph{#2}\par}
\newenvironment{knotetranscript}{\par\color{knotemuted}\small}{\par}
\newcommand\knotesummary[1]{\par\noindent\fcolorbox{knoteaccent}{knoteaccent!6}{%
  \parbox{\dimexpr\linewidth-2\fboxsep-2\fboxrule}{\textbf{Summary}\par #1}}\par}
\makeatletter
\newdimen\knote@limit \knote@limit=4950mm
\newbox\knote@page \newbox\knote@chunk
\def\knote@ht#1{\dimexpr\ht#1+\dp#1\relax}
\def\knote@setheight#1{\global\pdfpageheight=#1\relax
  \special{pdf:pagesize width \the\paperwidth\space height \the\pdfpageheight}}
\def\knote@flush{\ifvoid\knote@page\else
  \knote@setheight{\dimexpr\knote@ht\knote@page+40mm+2mm\relax}%
  \box\knote@page\par\clearpage\global\pdfpageheight=\paperheight\fi}
\newenvironment{knotechunk}
  {\par\setbox\knote@chunk=\vbox\bgroup\hsize=\textwidth\linewidth=\textwidth}
  {\par\egroup
   \ifdim\knote@ht\knote@chunk>\knote@limit
     \knote@flush\unvbox\knote@chunk\par\clearpage
   \else\ifvoid\knote@page
     \global\setbox\knote@page=\vbox{\unvbox\knote@chunk}%
   \else\ifdim\dimexpr\knote@ht\knote@page+\knote@ht\knote@chunk\relax>\knote@limit
     \knote@flush\global\setbox\knote@page=\vbox{\unvbox\knote@chunk}%
   \else
     \global\setbox\knote@page=\vbox{\unvbox\knote@page\unvbox\knote@chunk}%
   \fi\fi\fi}
\AtEndDocument{\knote@flush}
\makeatother
\pagestyle{empty}
\begin{document}
\begin{knotechunk}
{\color{knoteaccent}\rule{\linewidth}{1.5pt}}\par
{\LARGE\bfseries Seminar — Garnet solid electrolytes for lithium metal batteries}\par
{\large Speaker: Dr P. Varga}\hfill{\color{knotemuted}14 October 2026 \textperiodcentered{} Materials building, Seminar room}\par
{\color{knoteaccent}\rule{\linewidth}{0.6pt}}\par
\knotesummary{Garnet Li\textsubscript{7}La\textsubscript{3}Zr\textsubscript{2}O\textsubscript{12} (LLZO) is a promising solid electrolyte: stable against lithium metal and fast-conducting in its cubic form, which needs a dopant such as Al or Ga. Its weak point is the surface — Li\textsubscript{2}CO\textsubscript{3}/LiOH from air and chemical inhomogeneity at grain boundaries raise the interface resistance and let lithium penetrate.}

\section{Why solid electrolytes}

\begin{itemize}\setlength{\itemsep}{0pt}\setlength{\parskip}{0pt}
  \item replace the flammable liquid electrolyte
  \item enable a lithium metal anode: much higher energy density
  \item requirements: ionic conductivity \ensuremath{\geq} $10^{-4}$ S cm$^{-1}$, electronic insulation, stability against Li, good interfaces
\end{itemize}
\end{knotechunk}

\begin{knotechunk}
\section{Garnet LLZO}

$\mathrm{Li_7La_3Zr_2O_{12}}$ exists in two forms:

\begin{enumerate}\setlength{\itemsep}{0pt}\setlength{\parskip}{0pt}
  \item tetragonal (ordered Li): about $10^{-6}$ S cm$^{-1}$
  \item cubic (disordered Li): $10^{-4}$–$10^{-3}$ S cm$^{-1}$ at room temperature
\end{enumerate}

\knotekey{The cubic phase is stabilised by aliovalent doping — Al\textsuperscript{3+} or Ga\textsuperscript{3+} on Li sites, or Ta\textsuperscript{5+} / Nb\textsuperscript{5+} on Zr sites; Ga-doped LLZO is among the best conductors.}

\begin{itemize}\setlength{\itemsep}{0pt}\setlength{\parskip}{0pt}
  \item wide electrochemical window and kinetically stable against Li metal
\end{itemize}
\end{knotechunk}

\begin{knotechunk}
\section{The interface problem}

\begin{itemize}\setlength{\itemsep}{0pt}\setlength{\parskip}{0pt}
  \item in air, LLZO reacts with H\textsubscript{2}O and CO\textsubscript{2}: a LiOH / Li\textsubscript{2}CO\textsubscript{3} layer forms on the surface
  \item this layer is a poor Li\textsuperscript{+} conductor and is not wetted by lithium \ensuremath{\rightarrow} high interfacial resistance
  \item remedies: polishing, acid etching, heat treatment in inert gas, thin interlayers (e.g. Au, Al\textsubscript{2}O\textsubscript{3})
  \item dendrites: lithium can still penetrate, often along grain boundaries; the critical current density (CCD) is the measure
\end{itemize}

\subsection{Chemical inhomogeneity}

\begin{itemize}\setlength{\itemsep}{0pt}\setlength{\parskip}{0pt}
  \item dopants and impurities segregate to grain boundaries and the surface after processing
  \item surface-sensitive techniques see it: XPS (chemical state), ToF-SIMS (mapping, depth profiles), LEIS (low-energy ion scattering: the outermost atomic layer)
\end{itemize}
\end{knotechunk}

\begin{knotechunk}
\section{Measuring it}

\begin{itemize}\setlength{\itemsep}{0pt}\setlength{\parskip}{0pt}
  \item impedance spectroscopy: separate bulk, grain-boundary and interface arcs with an equivalent circuit
  \item galvanostatic cycling of Li | LLZO | Li cells at rising current to find the CCD
  \item stack pressure strongly changes the interface: report it
\end{itemize}

\knotequestion{How much of the improvement after surface treatment is chemistry, and how much is better contact?}
\end{knotechunk}
\end{document}
